\section{Giriş}
\label{sec:intro}

\subsection{Motivasyon: Silikon Büyümediğinde Modeli Küçültmek}

Makine öğrenmesinde son on yıla damgasını vuran yaklaşım oldukça
belirgindir: problem zorlaştıkça model büyütülür, daha güçlü bir
hızlandırıcı kullanılır ve bellek bütçesi genişletilir. Bu strateji
dil, görüntü ve çok kipli üretim alanlarında etkileyici sonuçlar
vermiştir~\cite{patterson2021carbon}. Bununla birlikte aynı
stratejinin enerji, karbon ve tedarik zinciri maliyeti de giderek
daha görünür hâle gelmiştir. 2020--2024 döneminde ileri seviye yarı
iletkenlerde yaşanan küresel arz sıkıntısı otomotiv, tüketici
elektroniği ve tıbbi cihaz sektörlerini
aksatmış~\cite{burkacky2022chipshortage,khan2021chipshortage};
sürekli çevrim içi çalışan çıkarım sistemlerinin enerji tüketimi ise
hem büyük ölçekli hem de hızla büyüyen bir yük hâline
gelmiştir~\cite{strubell2019energy,patterson2021carbon}.

Bu nedenle tamamlayıcı bir araştırma yönü giderek daha önemli hâle
gelmektedir. Buradaki soru, bir modeli çalıştırmak için hangi yeni
hızlandırıcıya ihtiyaç duyduğumuz değil; hâlihazırda kitlesel
ölçekte üretilen, her yıl on milyarlarca adet sevk edilen ve birim
maliyeti birkaç dolar düzeyinde kalan silikon üzerinde hangi
modellerin çalışabileceğidir. Bu silikon; giyilebilir cihazlarda,
sensörlerde, akıllı ev aygıtlarında, otomotiv alt sistemlerinde ve
endüstriyel telemetri uç noktalarında kullanılan 8-bit ve 16-bit
mikrodenetleyicilerdir (MCU). Bu alan genellikle ``tinyML''
olarak anılır~\cite{banbury2021mlperf}; GPU ölçekli çıkarımdan
farklı olarak belirli bir büyük ölçekli altyapıya bağlı değildir.

\subsection{Bare-Metal MCU Sınıfı}

Yayımlanmış tinyML uygulamalarının önemli bir bölümü ARM Cortex-M
sınıfı hedeflere yönelir. Bu aygıtlar yüzlerce kilobayt SRAM,
donanım kayan nokta birimleri, tek çevrimli çarpıcılar ve üretici
tarafından iyileştirilmiş sinir ağı çekirdekleri
sunar~\cite{lai2018cmsisnn,david2021tflitemicro,lin2020mcunet}.
Ucuzdurlar; ancak mikrodenetleyici dünyasının en alt ucuna göre
hâlâ oldukça güçlüdürler. Bu makalede incelediğimiz \emph{bare-metal}
MCU'lar ise kullanılabilir silikon yelpazesinin öteki ucunda
yer alır:
\begin{itemize}
\item 32~KB Flash, 2~KB SRAM, donanım $8{\times}8$ çarpıcısı ve
kayan nokta birimi bulunmayan 8-bit \arduino{} (ATmega328P);
\item 16~KB Flash, 512~B SRAM ve \emph{hiçbir tür donanım çarpıcısı
bulunmayan} 16-bit \msp{}. Bu aygıtta çarpma işlemleri yazılımla
öykünülür.
\end{itemize}
Bu iki parça bugün de aktif olarak üretilmektedir ve birim maliyetleri
ARM Cortex-M sınıfı hedeflerin belirgin biçimde altındadır. Bir
yinelemeli ağ \msp{} üzerinde gerçek zamanlı çalıştırılabiliyorsa,
aynı ağın bugün üretimde olan neredeyse her yaygın mikrodenetleyiciye
uyarlanabileceğini söylemek mümkündür.

\subsection{Neden Yinelemeli Ağlar, Neden \fastgrnn{}?}

Bileğe ya da bele takılan bir eylemsiz ölçüm sensöründen insan
aktivite tanıma (HAR), akış halinde sınıflandırma için doğal bir
örnektir. Giriş boyutu düşüktür (üç ivmeölçer kanalı), örnekleme
hızı sınırlıdır (50~Hz) ve çıkış uzayı küçüktür (altı sınıf)
\cite{anguita2013public,reyes2015transition}. Bu nedenle problem,
bare-metal MCU sınıfını değerlendirmek için uygundur: modele
sığacak kadar küçük, fakat pencerelenmiş ileri beslemeli bir taban
çizgisinden yarar görecek kadar zamansaldır.

\fastgrnn{}~\cite{kusupati2018fastgrnn}, kaynak kısıtlı çıkarım
için tasarlanmış kapılı bir yinelemeli hücredir. Düşük ranklı
ağırlık ayrıştırması, iteratif sert eşikleme (IHT) ile
seyrekleştirme~\cite{blumensath2009iht} ve Q15 sabit noktalı
nicelendirme~\cite{jacob2018integer,han2016deepcompression} gibi
birbirini tamamlayan sıkıştırma tekniklerini bir araya getirir.
Özgün çalışmada \fastgrnn{}'nun, LSTM~\cite{hochreiter1997lstm}
doğruluğuna çok daha az parametreyle yaklaşabildiği gösterilmiştir.
Ancak bu çalışma çarpıcısız bir hedef için herkese açık bare-metal
bir başvuru uygulaması sunmaz; ayrıca modelin akış kipindeki
davranışını ayrıntılı biçimde incelemez.

\subsection{Katkılar}

Bu makale, bare-metal MCU'larda HAR için \fastgrnn{} algoritmasının
uçtan uca ve açık kaynaklı bir yeniden üretimini sunar. Özgün
makalenin ötesinde dört katkı öne çıkmaktadır:

\begin{itemize}
\item \textbf{Platformlar arası deterministik çıkarım.} Tek bir
taşınabilir C kaynak dosyası hem 8-bit \arduino{} hem de 16-bit
\msp{} üzerinde derlenir. 3{,}399 test penceresinin tamamında
\emph{bit-eşdeğer} gizli durum yörüngeleri üretir ve PyTorch
referansıyla tahmin düzeyinde \%100 uyuşur. Donanımdan bağımsız bu
düzeyde yeniden üretilebilirlik, güvenlikle ilişkili gömülü
uygulamalar için özellikle değerlidir.

\item \textbf{Çarpıcısız hedefler için uygulanabilir bir arama
tablosu yöntemi.} $[-8, +8]$ giriş aralığında tanımlanan 256 girişli
sigmoid/tanh arama tablosu, \msp{} üzerinde tam pencere çıkarımını
\textbf{30.5$\times$} hızlandırır (54~s $\rightarrow$ 1.8~s), \%31
bütçe başı ile 50~Hz akış hedefini mümkün kılar. Bu yöntem
\fastgrnn{}'ya özgü değildir; $\sigma$ ya da $\tanh$ aktivasyonları
kullanan başka yinelemeli hücrelere de uygulanabilir.

\item \textbf{Yinelemeli ısınma gecikmesinin ölçülmesi.} Tahminlerin
kararlı hâle gelmesi için, donanımdan bağımsız olarak, yaklaşık
2~saniyelik ($\sim$100 örnek, 50~Hz) gizli durum geçmişi gerektiğini
gözlemliyoruz. Bu davranış donanımdan değil, yinelemeli dinamiklerden
kaynaklanır ve \fastgrnn{} sınıfı HAR sistemlerinde kullanıcıya
yansıyan yanıt süresini belirler. Özgün makalede bu nokta ele
alınmamıştır.

\item \textbf{Donanım enerji karakterizasyonu.} INA226 donanım akım
sensörü kullanılarak \msp{} hedef rayında ölçülen değerler:
\textbf{17.7~mW} aktif çıkarım gücü, \textbf{$<$0.09~mW} LPM3 uyku
gücü ve LUT etkinken 128 örneklik pencere başına \textbf{31.5~mJ}.
LUT olmadan aynı pencere 954~mJ tüketir; bu da arama tablosunun sabit
saat frekansında sağladığı 30.5$\times$ gecikme kazanımının doğrudan
sonucu olan \textbf{\%96.7 enerji azalmasına} karşılık gelir. Bu
karakterizasyon, modelin algoritmik verimliliği ile gerçek dünya güç
zarfı arasındaki bağı kurar ve özgün makalede bulunmayan bir ölçüm
boyutu ekler.
\end{itemize}

\noindent\textbf{Kod ve yeniden üretilebilirlik.} Kaynak kodu,
eğitilmiş modeller, dışa aktarılan Q15 ağırlıkları, deney günlükleri
ve gömülü hedefler için üretilen ikililer Apache Lisansı 2.0 altında
\url{https://github.com/emre1998/fastgrnn-har} adresinde
paylaşılmıştır. Tam yeniden üretim yaklaşık iki saat CPU eğitimi ve
beş dakika MCU aktarımı gerektirir.

\subsection{Makalenin Akışı}

Bölüm~\ref{sec:related}, kompakt RNN ve uç-ML literatürünü özetler.
Bölüm~\ref{sec:method}, \fastgrnn{} hücresini ve düşük ranklı,
seyrek, nicelendirilmiş (L-S-Q) sıkıştırma boru hattını tanımlar.
Bölüm~\ref{sec:setup} deneysel kurulumu açıklar. Bölüm~\ref{sec:results}
doğruluk, bellek ayak izi, gerçek zamanlı akış performansı ve donanım
enerji karakterizasyonunu raporlar. Bölüm~\ref{sec:discussion}, ısınma gecikmesini ve
platformlar arası determinizmi tartışır; ayrıca sınırlılıkları
verir. Bölüm~\ref{sec:conclusion} makaleyi sonuçlandırır.
